\section{Fundamentação e delimitação conceitual}
\subsection{GDD como especificação do objeto de análise}
No documento fornecido, o GDD organiza visão geral, estrutura do cenário, fluxo de missão, mecânicas, arquitetura, recursos visuais e backlog; seus enunciados têm níveis distintos de detalhe \gdd. O cenário de quatro pavimentos e o limiar de exposição à câmera oferecem critérios diretos de comparação, enquanto a referência ao terminal de hack não estabelece, no texto lido, sua duração nem todos os seus efeitos \gdd.

Essa heterogeneidade exige separar a presença de uma ideia de jogo da especificação suficiente de sua regra operacional: citar um temporizador não determina sua duração, e mencionar um retorno ao ponto inicial não resolve automaticamente a tolerância espacial \fonte. No estudo, a unidade de análise é um enunciado de decisão ou requisito associado a um critério verificável; um mesmo requisito amplo pode exigir mais de uma unidade para distinguir componentes e parâmetros \rep.

O antecedente de uso do ChatGPT na priorização de requisitos oferece um vínculo temático com decisões de engenharia de software, mas seus resultados não são transferidos para este jogo nem utilizados como validação da matriz proposta (Queiroz; Lima, 2025). A matriz é um instrumento operacional do caso, cuja clareza e consistência ainda dependem de aplicação e revisão humana documentadas \rep.

\subsection{Correspondência, ambiguidade e proveniência}
Correspondência designa a relação entre critério explícito e evidência examinada; ela não funciona, neste estudo, como índice global de qualidade, diversão ou valor pedagógico \rep. Uma regra pode divergir literalmente da especificação e ainda ser adequada para a interação, hipótese que só pode ser examinada com uma justificativa documentada ou verificação específica \fonte.

Ambiguidade e insuficiência de especificação são registradas separadamente da categoria de correspondência, porque a ausência de parâmetro numérico não equivale à comprovação de uma implementação incorreta \rep. A categoria não verificável conserva essa insuficiência como resultado documental, em vez de converter silenciosamente a escolha do código em requisito do GDD \rep.

Proveniência da decisão é a identificação do registro que sustenta quem resolveu uma ambiguidade ou definiu um parâmetro; o relato geral de programação por agente não identifica o autor de cada escolha \rep. Na ausência de conversas originais, solicitações ou decisões assinadas, a matriz conserva o responsável como desconhecido, inclusive nos casos em que o código contém um valor numérico claro \rep.

\subsection{Artefato técnico e evidência educacional}
Prather et al. (2024) examinam benefícios e danos da IA generativa para programadores iniciantes, oferecendo fundamento para cautela ao associar desempenho em tarefas assistidas a competência individual. Por isso, a existência do jogo e do GDD é tratada aqui como evidência documental da experiência relatada, sem medida de aprendizagem, comparação entre estudantes ou conclusão sobre eficácia pedagógica (Prather et al., 2024; Night Heist, 2026).

Lelli, Santos e Castelo Branco (2024) apresentam práticas gamificadas no ensino de verificação e validação, fornecendo um antecedente educacional próximo ao interesse pelo desenvolvimento de jogos. O vínculo temático não torna o Night Heist uma intervenção educacional avaliada: o corpus atual sustenta a análise técnica de decisões e não resultados de ensino (Lelli; Santos; Castelo Branco, 2024; Night Heist, 2026).
